% Latex template: mahmoud.s.fahmy@students.kasralainy.edu.eg
% For more details: https://www.sharelatex.com/learn/Beamer

\documentclass{beamer}					% Document class

\usepackage[english]{babel}				% Set language
\usepackage[utf8x]{inputenc}			% Set encoding

\mode<presentation>						% Set options
{
  \usetheme{default}					% Set theme
  \usecolortheme{default} 				% Set colors
  \usefonttheme{default}  				% Set font theme
  \setbeamertemplate{caption}[numbered]	% Set caption to be numbered
}

\usepackage{graphicx}					% For including figures
\usepackage{booktabs}					% For table rules
\usepackage{hyperref}					% For cross-referencing
\usepackage{subcaption}
\usepackage{multicol}
\usepackage{listings}					% For code listings
\usepackage{xcolor}

% --- R code style for listings ---
\definecolor{codegreen}{rgb}{0.0,0.5,0.2}
\definecolor{codegray}{rgb}{0.5,0.5,0.5}
\definecolor{codeblue}{rgb}{0.1,0.2,0.7}
\definecolor{backcolour}{rgb}{0.96,0.96,0.94}

\lstdefinestyle{Rstyle}{
    backgroundcolor=\color{backcolour},
    commentstyle=\color{codegreen},
    keywordstyle=\color{codeblue}\bfseries,
    stringstyle=\color{codegreen},
    basicstyle=\ttfamily\footnotesize,
    breaklines=true,
    captionpos=b,
    keepspaces=true,
    showspaces=false,
    showstringspaces=false,
    showtabs=false,
    tabsize=2,
    language=R,
    morekeywords={chisq.test, matrix, c, fisher.test, mcnemar.test, table}
}
\lstset{style=Rstyle}

\title{The Chi-squared Test}	% Presentation title
\author{Rex Shao-Yu Wang}								% Presentation author
\institute{Department of Agronomy, National Taiwan University}					% Author affiliation
\date{June 1, 2026}									% Today's date	

\begin{document}

% Title page
\begin{frame}
  \titlepage
\end{frame}

% Outline
\begin{frame}{Outline}
  \tableofcontents
\end{frame}

%==================================================================
\section{Quick Review}
%==================================================================

\begin{frame}{When Do We Use the Chi-squared Test?}

\small
\begin{block}{The Decision Rule}
Both the response (Y) \textbf{and} the predictor (X) are \textbf{categorical}
(counts / frequencies, not means).
\end{block}

\vspace{0.6em}
\begin{center}
\begin{tabular}{ll}
\toprule
\textbf{Data type} & \textbf{Test} \\
\midrule
Continuous Y, continuous X & Correlation / Regression \\
Continuous Y, categorical X & t-test / ANOVA \\
\textbf{Categorical Y, categorical X} & \textbf{Chi-squared test} \\
\bottomrule
\end{tabular}
\end{center}

\vspace{0.8em}
\begin{alertblock}{Core Idea}
The core idea is always: compare \textbf{Observed} vs \textbf{Expected} counts.
$$\chi^2 = \sum \frac{(O - E)^2}{E}$$
\end{alertblock}

\end{frame}


\begin{frame}{Three Chi-squared Tests --- One Function in R}
\begin{itemize}
    \item \textbf{Test of Independence}\\
    \hspace{1em} Tests whether {two categorical variables are associated}
    
    \vspace{0.6em}
    \item \textbf{Goodness-of-Fit Test} \\
    \hspace{1em} Tests whether {observed frequencies match a specified distribution}
    
    \vspace{0.6em}
    \item \textbf{Test of Homogeneity} \\
    \hspace{1em} Tests whether {multiple populations share the same distribution}
\end{itemize}
\small
\begin{center}
\begin{tabular}{lll}
\toprule
\textbf{Test} & \textbf{Question} & \textbf{R function} \\
\midrule
Goodness-of-Fit & Fit a theoretical ratio? & \texttt{chisq.test(x, p=)} \\
Independence & Two variables associated? & \texttt{chisq.test(table)} \\
Homogeneity & Same distribution across groups? & \texttt{chisq.test(table)} \\
\bottomrule
\end{tabular}
\end{center}


\end{frame}

%==================================================================
\section{Building the Data}
%==================================================================

\begin{frame}[fragile]{Two Ways to Build a Contingency Table}

\textbf{Option A --- Type the counts directly with \texttt{matrix()}}
\begin{lstlisting}
oil <- matrix(c(154, 132,
                180, 126,
                104, 131),
              nrow = 3, byrow = TRUE)
rownames(oil) <- c("Support", "Oppose", "DoNotKnow")
colnames(oil) <- c("Grad", "NonGrad")
\end{lstlisting}

\vspace{0.4em}
\textbf{Option B --- From raw data with \texttt{table()}}
\begin{lstlisting}
# df has one row per respondent
oil <- table(df$opinion, df$college)
\end{lstlisting}
\end{frame}

%==================================================================
\section{R Implementation}
%==================================================================

\begin{frame}{Goodness-of-Fit Test}
\small
\begin{block}{What it tests}
Whether the observed counts of a \textbf{single categorical variable} match a
\textbf{theoretical / expected ratio}.
\end{block}

\vspace{0.5em}
$$
\begin{cases}
H_0: \text{the population follows the given distribution} \\
H_1: \text{it does not}
\end{cases}
$$

\begin{itemize}
    \item Expected count: $E_i = p_i \times N$ (a known ratio, no second variable)
    \item df $= k - 1$, where $k$ = number of categories
\end{itemize}

\vspace{0.8em}
\begin{alertblock}{Keep in mind}
Here the expected proportions come from \textbf{theory} (e.g.\ a genetic law),
not from the data itself.
\end{alertblock}

\end{frame}


\begin{frame}[fragile]{Goodness-of-Fit Test}
\small
\textbf{Example:} Do budgie colours follow Mendel's \textbf{9:3:3:1} ratio? \\
Counts of $N = 971$ birds: Green 550, Blue 168, Yellow 180, White 73. ($\alpha=0.05$)

\begin{lstlisting}
observed <- c(Green=550, Blue=168, Yellow=180, White=73)
expected_p <- c(9, 3, 3, 1) / 16  

chisq.test(x = observed, p = expected_p)
\end{lstlisting}

\vspace{0.3em}
\texttt{X-squared = 3.6342, df = 3, p-value = 0.3038} \\
\vspace{0.3em}
$\because$ p $>$ 0.05 \ $\therefore$ data are consistent with Mendel's laws.

\end{frame}


\begin{frame}{Test of Independence}
\small
\begin{block}{What it tests}
Whether \textbf{two categorical variables} measured on \textbf{one sample}
are associated (related) or independent.
\end{block}

\vspace{0.5em}
$$
\begin{cases}
H_0: \text{the two factors are independent} \\
H_1: \text{ the two factors are associated}
\end{cases}
$$
\begin{itemize}
    \item Data are arranged in a \textbf{contingency table}
    \item Expected count: $E_{ij} = (R_i \times C_j) / N$ --- estimated \textbf{from the data}
    \item df $= (r-1)(c-1)$
\end{itemize}

\vspace{0.8em}
\begin{alertblock}{Keep in mind}
Unlike goodness-of-fit, there is \textbf{no theoretical distribution} ---
expected counts come from the row and column totals.
\end{alertblock}

\end{frame}


\begin{frame}[fragile]{Test of Independence}
\small
\textbf{Example:} 827 people were asked their opinion on oil drilling,
split by college status. Is opinion \textbf{associated} with having a degree?

\begin{lstlisting}
oil <- matrix(c(154, 132,
                180, 126,
                104, 131),
              nrow = 3, byrow = TRUE)

chisq.test(oil, correct = FALSE) 
chisq.test(oil)$expected 
\end{lstlisting}

\vspace{0.3em}
\texttt{X-squared = 11.47, df = 2, p-value = 0.0032}\\
\vspace{0.3em}
$\because$ p $<$ 0.05 \ $\therefore$ reject $H_0$: the two factors are NOT independent.

\vspace{0.6em}
\begin{alertblock}{Reminder}
For a \textbf{2x2} table, \texttt{chisq.test()} applies Yates' correction by
\textbf{default}. Use \texttt{correct = FALSE} to turn it off.
\end{alertblock}
\end{frame}


\begin{frame}{Test of Homogeneity}
\small
\begin{block}{What it tests}
Whether \textbf{two (or more) separate populations} share the
\textbf{same distribution} over the categories of one variable.
\end{block}

\vspace{0.5em}
$$
\begin{cases}
H_0: \text{the populations follow the same distribution} \\
H_1: \text{at least one population differs}
\end{cases}
$$

\begin{itemize}
    \item Same contingency table, same $E_{ij} = (R_i \times C_j)/N$, same df
\end{itemize}

\vspace{0.8em}
\begin{alertblock}{Independence vs Homogeneity}
The \textbf{computation is identical} to the independence test.\\
The difference is \textbf{study design}: 
\begin{itemize}
    \item Independence: \textbf{one} sample cross-classified
    \item Homogeneity: \textbf{separate} samples from each population
\end{itemize}
\end{alertblock}

\end{frame}


\begin{frame}[fragile]{Test of Homogeneity}
\small
\textbf{Example:} Separate samples of male and female students were asked
their living arrangement. Do the two groups share the \textbf{same distribution}?

\begin{lstlisting}
living <- matrix(c(72, 84, 49, 45,
                   91, 86, 88, 35),
                 nrow = 2, byrow = TRUE)
rownames(living) <- c("Male", "Female")
colnames(living) <- c("Dorm","Apt","Home","Other")

chisq.test(living)
\end{lstlisting}

\vspace{0.3em}
\texttt{X-squared = 10.13, df = 3, p-value = 0.0175}\\

\vspace{0.3em}
$\because$ p $<$ 0.05 \ $\therefore$ reject $H_0$: the distributions are NOT the same.

\end{frame}


\begin{frame}{Interpreting the R Output}
\small
A typical \texttt{chisq.test()} output gives three numbers:

\vspace{0.6em}
\begin{itemize}
    \item \textbf{X-squared} --- the test statistic $\chi^2_0$
    \item \textbf{df} --- degrees of freedom
    \item \textbf{p-value} --- probability of this result if $H_0$ were true
\end{itemize}

\vspace{0.8em}
\begin{block}{Decision Rule}
If \textbf{p-value $< \alpha$} (significance level) $\Rightarrow$ \textbf{reject $H_0$}.
\end{block}

\end{frame}

%==================================================================
\section{Paired Data: McNemar Test}
%==================================================================
 
\begin{frame}{McNemar Test}
\small
\begin{block}{When to use}
When the two samples are \textbf{paired / not independent} --- the
\textbf{same subjects} measured twice (e.g.\ the same voters before vs
after a debate).
\end{block}
 
\vspace{0.5em}
\begin{itemize}
    \item $H_0$: the proportion is unchanged (no shift between the two times)
    \item Only the \textbf{off-diagonal} cells ($n_{12}$, $n_{21}$) carry information ---
          the agreements on the diagonal cancel out
    \item Test statistic: $\chi^2 = \dfrac{(n_{12} - n_{21})^2}{n_{12} + n_{21}}$, df $= 1$
\end{itemize}
 
\vspace{0.8em}
\begin{alertblock}{Why not the ordinary test?}
The ordinary chi-squared test assumes \textbf{independent} observations.
Paired data violate this, so a different statistic is needed.
\end{alertblock}
 
\end{frame}
 
 
\begin{frame}[fragile]{McNemar Test}
\small
\textbf{Example:} 1000 voters stated their choice \textbf{before} and
\textbf{after} a debate. Did the debate change their decision?
 
\vspace{0.4em}
\centering
\begin{tabular}{lcc}
\toprule
 & After: A & After: B \\
\midrule
Before: A & 491 & 9 \\
Before: B & 1 & 499 \\
\bottomrule
\end{tabular}
 
\flushleft
\vspace{0.5em}
\begin{lstlisting}
votes <- matrix(c(491, 9,
                  1, 499),
                nrow = 2, byrow = TRUE)
 
mcnemar.test(votes, correct = FALSE)
\end{lstlisting}
 
\vspace{0.3em}
\texttt{McNemar's chi-squared = 6.4, df = 1, p-value = 0.0114}\\
\vspace{0.3em}
$\because$ p $<$ 0.05 \ $\therefore$ reject $H_0$: the debate did change voters' decisions.
 
\end{frame}
 
%==================================================================
\section{Summary}
%==================================================================
 
\begin{frame}[fragile]{Summary: The R Workflow}
\small
\begin{enumerate}
    \item Build the table --- \texttt{matrix()} or \texttt{table()}
    \item Run the test --- \texttt{chisq.test()}
    \item \textbf{Check} expected counts --- \texttt{result\$expected}
    \item Read the p-value and decide
\end{enumerate}
 
\vspace{1em}
\begin{center}
\begin{tabular}{ll}
\toprule
\textbf{Situation} & \textbf{Function} \\
\midrule
Fit a ratio & \texttt{chisq.test(x, p=)} \\
Independence / Homogeneity & \texttt{chisq.test(table)} \\
Paired data & \texttt{mcnemar.test()} \\
\bottomrule
\end{tabular}
\end{center}
 
\vspace{0.8em}
\centering \textit{Same idea every time: Observed vs Expected.}
 
\end{frame}
 
\begin{frame}{}    
    \begin{center}
    \Large Thank you for your attention!
    \end{center}
\end{frame}
 
\end{document}